% Pista ana-24j-q3 · Anàlisi
% Procedència: PAU Catalunya 2024 · Convocatòria de juny · Sèrie 1 · Qüestió 3
\documentclass[11pt,a4paper]{article}
\usepackage{pista}
\pistainfo{Catalunya 2024}{Convocatòria de juny}{Sèrie 1}

\begin{document}

\section*{\textcolor{blauPista}{Qüestió 3}}

\noindent En Joan troba entre els papers del seu avi un esbós com el de la figura adjunta, on es descriu un terreny de regadiu que ha deixat en herència al seu pare. La corba de la gràfica és $y = f(x)$, amb $f(x) = -x^{3} + 7x^{2} - 6x + 5$.

\subsection*{\textit{a)} \normalfont A partir de l'expressió de $f(x)$, calculeu les coordenades dels punts $P$, $Q$ i $R$ indicats a la figura. Calculeu també l'equació de la recta $PR$. \hfill\textnormal{[1,25 punts]}}

\pista{Observa la figura amb cura: $P$ i $Q$ estan sobre la recta horitzontal $y = 5$, mentre que $R$ es troba sobre l'eix $OX$ (és a dir, a alçada $y = 0$) i sembla que té la mateixa abscissa que $Q$. Per a trobar $P$ i $Q$, resol l'equació $f(x) = 5$. Per trobar el punt $R$, n'hi ha prou amb llegir la figura un cop tinguis $Q$. Amb $P$ i $R$, la recta $PR$ queda determinada i pots trobar la seva equació, de forma senzilla, com a punt-pendent.}

\subsection*{\textit{b)} \normalfont Calculeu la superfície del terreny. \hfill\textnormal{[1,25 punts]}}

\pista{La regió ombrejada està delimitada per dalt per la corba $y = f(x)$ i per baix per la recta $PR$. L'àrea és la integral definida de la diferència entre les dues funcions (la de ``dalt'' menys la de ``baix''), entre les abscisses dels dos punts de tall corresponents (que ja tens de l'apartat a).}

\end{document}
